% Annexe F -- Exemples de code : Python, Rust (rusty-SUNDIALS) et LeanFlow
% A placer apres \appendix. Chaque listing est lu tel quel (\lstinputlisting) dans code/examples/, et chaque
% sortie dans code/examples/logs/ : ce sont les fichiers executes le 2026-10-26, rien n'est recopie a la main.
\chapter{Exemples de code : Python, Rust (rusty-SUNDIALS) et LeanFlow}\chaptermark{Exemples de code}\label{appF:chap}

\lstdefinelanguage{Rust}{morekeywords={fn,let,mut,pub,use,struct,impl,for,in,if,else,match,return,move,loop,
  const,type,as,mod,where,true,false,Some,None,Ok,Err,Vec,String,f64,usize,u64,Result,Self,self},
  sensitive=true,morecomment=[l]{//},morestring=[b]"}
\lstdefinestyle{appFcode}{basicstyle=\ttfamily\scriptsize,numbers=left,numberstyle=\tiny\color{gray},
  numbersep=4pt,xleftmargin=1.2em}
\lstdefinestyle{appFlog}{language={},basicstyle=\ttfamily\scriptsize,backgroundcolor=\color{gray!6},
  frame=single,rulecolor=\color{tierX!40},
  literate={/}{{/\allowbreak}}1 {ℕ}{{$\mathbb{N}$}}1 {ℝ}{{$\mathbb{R}$}}1 {→}{{$\to$}}1 {≠}{{$\ne$}}1
  {∀}{{$\forall$}}1 {⁻¹}{{$^{-1}$}}1}
\newcommand{\appFsortie}[1]{\par\noindent\textbf{Sortie réelle} (\texttt{\detokenize{#1}}) :%
  \lstinputlisting[style=appFlog]{code/examples/logs/#1}}

Les chapitres~1 à~5 énoncent des théorèmes et racontent comment ils ont été vérifiés. Cette annexe donne au
lecteur les outils pour \emph{refaire} les calculs qui les accompagnent. Elle contient douze exemples, répartis sur les
thèmes du livre. Chacun suit le même plan :
\begin{enumerate}
\item \textbf{la question}, posée avant tout calcul ;
\item \textbf{le code}, un extrait court lu directement dans le fichier exécuté (les numéros de ligne sont ceux du
fichier) ;
\item \textbf{la sortie réelle} du programme, recopiée par \LaTeX{} depuis le journal de l'exécution ;
\item \textbf{une figure}, quand elle aide à lire les nombres ;
\item \textbf{un exercice}.
\end{enumerate}

\begin{warnbox}[: ce que ces exemples sont, et ce qu'ils ne sont pas]
\begin{itemize}
\item Un calcul en virgule flottante est de niveau~X : il illustre, il teste, il peut \emph{réfuter}, mais il ne
démontre jamais. Seuls les exemples~\ref{appF:ex10} et~\ref{appF:ex11} touchent au noyau de Lean.
\item Chaque programme exécute au moins un témoin positif (un cas dont on connaît la réponse) et un témoin négatif (un
cas faux exprès, qui doit être refusé). Il se termine avec le code de sortie~1 si un témoin ne se comporte pas comme
prévu. Toutes les sorties imprimées ici se terminent par \texttt{temoins conformes: True} ou
\texttt{controls as expected: true}, sauf celle de l'exemple~\ref{appF:ex11}. Celle-ci est la sortie brute de Lean,
qui compile avec le code~0 même en présence du témoin planté ; le témoin se lit alors dans \texttt{\#print axioms}.
\item Le code, les sorties et les figures de cette annexe ont été produits par un agent d'IA (Claude) sous la direction
de l'auteur, et vérifiés seulement par cet agent. Aucun être humain ne les a audités.
\end{itemize}
\end{warnbox}

\paragraph{Où sont les fichiers.} Tout est rangé dans \texttt{code/examples/} :
\begin{itemize}
\item \texttt{python/} : les exemples Python, et \texttt{appF\_style.py}, le style commun des figures ;
\item \texttt{rust/} : une copie octet par octet de la crate \texttt{qf-textbook-examples} de rusty-SUNDIALS :
\begin{itemize}
\item branche \texttt{textbook2/code-examples-2026-10-26}, issue de \texttt{main} au commit \texttt{7630f7e} ;
\item commit \texttt{f01b82b730c86ccf05c1f773b11fe0a9b650cc7b} ;
\end{itemize}
\item \texttt{lean/} : le fichier Lean de l'exemple~\ref{appF:ex11} et les deux certificats engendrés par LeanFlow à
l'exemple~\ref{appF:ex10} ;
\item \texttt{logs/} : les sorties imprimées ici ; les fichiers CSV écrits par les binaires Rust sont dans
\texttt{logs/rust\_out/}.
\end{itemize}

\paragraph{Reproduire.} Les exemples Python utilisent numpy, scipy, mpmath et matplotlib. On les lance avec
l'interpréteur isolé (option \texttt{-I}) :
\begin{lstlisting}[language={}]
python -I code/examples/python/f01_poisson.py > code/examples/logs/f01_poisson.txt
\end{lstlisting}
La crate Rust se construit dans la copie de travail de rusty-SUNDIALS :
\begin{lstlisting}[language={}]
CARGO_TARGET_DIR=~/rusty-target-ar cargo build --release -j 2 --offline -p qf-textbook-examples
~/rusty-target-ar/release/tb_moran_forward code/examples/logs/rust_out
\end{lstlisting}
Avant la validation, nous avons passé trois des contrôles de l'intégration continue du dépôt, et tous trois ont
réussi :
\begin{itemize}
\item \texttt{cargo fmt --all -- --check} ;
\item \texttt{cargo clippy --workspace --features experimental -- -D warnings} ;
\item \texttt{cargo test -p qf-textbook-examples}, avec 5 tests sur 5 réussis.
\end{itemize}
La suite de tests complète de l'espace de travail n'a pas été relancée.

\begin{table}[ht]\centering\small
\caption{Les douze exemples. Niveau : celui de l'énoncé \emph{vérifié par le programme}, pas celui du théorème qu'il
illustre.}\label{appF:tab}
\begin{tabular}{llp{6.1cm}l}\toprule
& langage & question & niveau\\\midrule
1 & Python & formule de Poisson pour les gaussiennes & X\\
2 & Python & énergie des réseaux triangulaire et carré (gaussien, Yukawa, dipolaire) & X\\
3 & Python & signes de $(-1)^ng^{(n)}$ avec mpmath, et deux contre-exemples & X\\
4 & Python & puissances exactes de Wright--Fisher, $P^n(i,N)\to i/N$ & exact (fractions)\\
5 & Python & fixation sur un graphe isotherme et sur l'étoile & exact (fractions)\\
6 & Python & phonons et énergie de point zéro, maximum local & X\\
7 & Rust & CVODE sur une EDO raide, puis l'équation de Kolmogorov de Moran & X\\
8 & Rust & relaxation d'un petit cristal de Yukawa & X\\
9 & Rust & les invariants du modèle en couches $\sigma=3$ & X\\
10 & LeanFlow & certifier des fractions de fixation, rejeter une affirmation plantée & A (arithmétique)\\
11 & Lean & lire et vérifier \lean{ILENE} & A\\
12 & Python & contrôle croisé CVODE contre SciPy & X\\\bottomrule
\end{tabular}
\end{table}

% =====================================================================================================
\section{Exemple 1 -- La formule de Poisson pour les gaussiennes}\label{appF:ex1}

\paragraph{La question.} La formule de sommation de Poisson est le moteur des chapitres~1 et~3. Pour une gaussienne, en
dimension~1, elle s'écrit
\[\sum_{n\in\Z}e^{-\pi a(n+x)^2}=a^{-1/2}\sum_{k\in\Z}e^{-\pi k^2/a}\cos(2\pi kx),\qquad a>0 .\]
Pour le réseau triangulaire $\Lat$ de covolume~1, dont le dual est une rotation de lui-même, elle donne
$\Theta_\Lat(a)=a^{-1}\Theta_\Lat(1/a)$. Peut-on la voir à $10^{-15}$ près ? Et à quoi sert-elle en pratique ?

\paragraph{Le code.} Les deux côtés de la formule, et la fonction thêta du réseau triangulaire
($|v|^2=\tfrac{2}{\sqrt3}(m^2+mn+n^2)$) :
\lstinputlisting[language=Python,style=appFcode,firstline=23,lastline=38,firstnumber=23]{code/examples/python/f01_poisson.py}
\appFsortie{f01_poisson.txt}

\paragraph{Lecture.} Les deux côtés coïncident à la précision de la machine ($10^{-16}$). Sans le facteur $a^{-1/2}$
(ou $a^{-1}$ en dimension~2), l'écart est d'ordre~1 : c'est le témoin négatif. La dernière ligne montre l'intérêt
pratique de la formule. Pour $a=0{,}01$, la gaussienne est très large : la somme directe demande $n_{\max}=28$ pour
atteindre $10^{-12}$, alors que le côté dual est exact dès son premier terme. Ce principe, sommer du côté où la
fonction décroît vite, est celui de la sommation d'Ewald.

\begin{figure}[ht]\centering
\includegraphics[width=\textwidth]{appF_poisson.pdf}
\caption{Exemple~1 \tierX. (a) $\Theta_\Lat(a)$ et $a^{-1}\Theta_\Lat(1/a)$ coïncident ; sans le facteur $a^{-1}$ (tirets),
la courbe est fausse. (b) Erreur relative de la somme directe en fonction du nombre de termes, pour $a=0{,}01$, $x=0{,}3$.}
\label{appF:fig:poisson}
\end{figure}

\begin{exercise}[Exemple 1]\label{appF:exo1}
(a) Démontrer la formule de dimension~1 à partir de la transformée de Fourier de $e^{-\pi ax^2}$, qui vaut
$a^{-1/2}e^{-\pi\xi^2/a}$. (b) Modifier \texttt{theta1} pour une autre fonction paire, $e^{-|x|}$, dont la
transformée est $2/(1+4\pi^2\xi^2)$, et vérifier numériquement la formule correspondante. Pourquoi faut-il beaucoup
plus de termes du côté dual ?
\end{exercise}

% =====================================================================================================
\section{Exemple 2 -- Réseau triangulaire contre réseau carré}\label{appF:ex2}

\paragraph{La question.} À densité~1, l'énergie par particule $E(L)=\tfrac12\sum_{v\in L\setminus0}V(|v|)$ est-elle plus
basse pour le réseau triangulaire que pour le réseau carré ? On prend trois potentiels complètement monotones en $r^2$ :
\begin{itemize}
\item le gaussien $e^{-\pi r^2}$, pour lequel $E=(\Theta_L(1)-1)/2$ ;
\item le Yukawa $e^{-\kappa r}/r$, pour $\kappa=1$ et $\kappa=2$ ;
\item le dipolaire $V=2(1/r-1/\sqrt{r^2+d^2})$, pour $d=3$, dans les unités $U=V/d^2$ de la crate
\texttt{qf-yukawa-crystal}.
\end{itemize}
Pour ces trois potentiels, l'optimalité du réseau triangulaire est \tierA{} (chapitre~2), conditionnelle au théorème
amont \lean{universal_energy_minimum}.

\paragraph{Le code.} Le potentiel dipolaire décroît comme $1/r^3$. Une coupure brutale à $|v|<R$ laisse une erreur de
l'ordre de $10^{-5}$, même en ajoutant la queue continue (la sortie le montre). On coupe donc en douceur, avec un
poids $w=\tfrac12\operatorname{erfc}((r-R)/\sigma)$, et l'on intègre le reste $(1-w)U$ comme un continuum de densité~1 :
\lstinputlisting[language=Python,style=appFcode,firstline=46,lastline=64,firstnumber=46]{code/examples/python/f02_lattice_energies.py}
\appFsortie{f02_lattice_energies.txt}

\paragraph{Lecture.} Dans les trois cas, le carré coûte plus cher. La dernière colonne compare aux valeurs
enregistrées par rusty-SUNDIALS :
\begin{itemize}
\item pour Yukawa, aux sommes exactes de \texttt{exact\_refs.jsonl}, avec des écarts de $10^{-16}$ ;
\item pour le dipolaire, aux énergies d'Ewald de \texttt{lattice\_gaps.jsonl}, avec un écart de $5\times10^{-13}$.
\end{itemize}
C'est le témoin positif. Le témoin négatif rappelle une règle des énoncés : à \emph{pas} égal, le réseau triangulaire
est plus dense, son énergie est plus haute et la comparaison s'inverse. Elle n'a de sens qu'à densité fixée.
Sur la figure~\ref{appF:fig:energies}a, les trois courbes normalisées sont presque confondues et croissent sans
interruption de $60^\circ$ à $90^\circ$. Ce n'est qu'un chemin parmi d'autres dans l'espace des réseaux : le théorème
porte sur toutes les configurations.

\begin{figure}[ht]\centering
\includegraphics[width=\textwidth]{appF_lattice_energies.pdf}
\caption{Exemple~2 \tierX. (a) Énergie le long des réseaux de forme $\tau=e^{i\theta}$, normalisée entre le triangulaire
($\theta=60^\circ$) et le carré ($\theta=90^\circ$), pour les trois potentiels. (b) Les deux réseaux de covolume~1.}
\label{appF:fig:energies}
\end{figure}

\begin{exercise}[Exemple 2]\label{appF:exo2}
(a) Remplacer \texttt{"gauss"} par une gaussienne $e^{-\pi ar^2}$ de paramètre $a$ quelconque, et vérifier sur une
grille de $a$ que $\Theta_\Lat(a)\le\Theta_{\Z^2}(a)$ (théorème de Montgomery, \tierL). (b) Pour un potentiel
\emph{non} complètement monotone, par exemple $V=e^{-r^2}-\tfrac12e^{-r^2/4}$, tracer l'énergie le long du même arc.
Le minimum est-il toujours à $60^\circ$ ? Que peut-on en conclure, et que ne peut-on pas en conclure ?
\end{exercise}

% =====================================================================================================
\section[Exemple 3 -- Complète monotonie]{Exemple 3 -- Complète monotonie : ce qu'un ordinateur peut dire}\label{appF:ex3}

\paragraph{La question.} Une fonction $g$ est complètement monotone (CM) si $(-1)^ng^{(n)}(t)\ge0$ pour tout $n$ et tout
$t>0$. Un programme ne teste qu'un nombre fini de couples $(n,t)$. Il peut donc \emph{réfuter} la complète monotonie,
puisqu'un seul signe faux suffit, mais jamais la démontrer : la preuve appartient à la bibliothèque CM/Bernstein et
au théorème \lean{ALEXIS} du chapitre~5. Que voit-on sur trois potentiels CM et sur deux fonctions qui ne le sont
pas ?

\paragraph{Le code.} \texttt{mpmath.diff} dérive numériquement à 40 chiffres significatifs, assez pour l'ordre~8 :
\lstinputlisting[language=Python,style=appFcode,firstline=22,lastline=40,firstnumber=22]{code/examples/python/f03_cm_signs.py}
\appFsortie{f03_cm_signs.txt}

\paragraph{Lecture.} Les trois potentiels physiques passent les $9\times7$ tests. Les deux contre-exemples échouent, et
la sortie dit où :
\begin{itemize}
\item $e^{-t}-\tfrac12e^{-2t}$ est un mélange d'exponentielles dont la mesure de Bernstein $\delta_1-\tfrac12\delta_2$
change de signe. Il échoue dès $n=2$, car $(-1)^ng^{(n)}=e^{-t}-2^{n-1}e^{-2t}<0$ pour $t$ petit.
\item Lennard-Jones échoue dès $n=0$ : il devient négatif.
\end{itemize}
Le premier échec est un certificat de non-complète monotonie, que l'on peut vérifier à la main. Les succès, eux, ne
certifient rien.

\begin{figure}[ht]\centering
\includegraphics[width=\textwidth]{appF_cm_signs.pdf}
\caption{Exemple~3 \tierX. Signe de $(-1)^ng^{(n)}(t)$ (bleu : $+$, orange : $-$, blanc : zéro exact) pour $n\le8$ et sept
valeurs de $t$. Les trois premiers potentiels sont CM (théorèmes \tierA{} du chapitre~2) ; les deux derniers ne le
sont pas.}\label{appF:fig:cm}
\end{figure}

\begin{exercise}[Exemple 3]\label{appF:exo3}
(a) Démontrer à la main que $e^{-t}-\tfrac12e^{-2t}$ n'est pas CM, en partant de la valeur de $g''(0^+)$. (b) Le
programme dit que $g(t)=e^{-\sqrt t}$ passe tous les tests. Démontrer qu'elle est CM en l'écrivant comme
$f\circ\phi$, avec $f=e^{-s}$ CM et $\phi=\sqrt{\cdot}$ de Bernstein (chapitre~1). (c) Ajouter à \texttt{FUNCS}
la fonction $1/(1+t^2)$ et prédire, avant de lancer le programme, l'ordre $n$ du premier échec.
\end{exercise}

% =====================================================================================================
\section{Exemple 4 -- Wright--Fisher en fractions exactes}\label{appF:ex4}

\paragraph{La question.} Dans le modèle de Wright--Fisher neutre à $N$ copies d'un gène, la génération suivante est un
tirage binomial : $P(i,j)=\binom Nj(i/N)^j(1-i/N)^{N-j}$. Le théorème \lean{KAL} (absorption de Kolmogorov,
\tierA) dit que $P^n(i,N)\to h(i)$ pour toute fonction $h$ harmonique avec $h(0)=0$ et $h(N)=1$ (et une condition de
variance). Ici $h(i)=i/N$. Peut-on le voir \emph{exactement}, sans arrondi ? À quelle vitesse la convergence a-t-elle
lieu ?

\paragraph{Le code.} Le module \texttt{fractions} de Python calcule en rationnels exacts :
\lstinputlisting[language=Python,style=appFcode,firstline=22,lastline=32,firstnumber=22]{code/examples/python/f04_wright_fisher.py}
\appFsortie{f04_wright_fisher.txt}

\paragraph{Lecture.} Deux égalités sont \emph{exactes}, et non à $10^{-16}$ près : la matrice est stochastique et
$Ph=h$. Avec sélection ($r=2$), $h(i)=i/N$ n'est plus harmonique : c'est le témoin négatif. Les dénominateurs
croissent vite, avec 152 chiffres à $n=60$. Le rapport de deux écarts successifs vaut $0{,}833333333333=1-1/N$,
c'est-à-dire la deuxième valeur propre de la matrice de Wright--Fisher, $\lambda_2=1-1/N$.

\begin{figure}[ht]\centering
\includegraphics[width=\textwidth]{appF_wright_fisher.pdf}
\caption{Exemple~4. (a) $P^n(i,N)$ pour $n=1$, $3$ et $20$, et la limite $i/N$, pour $N=6$. (b) L'écart maximal, calculé
exactement puis converti en flottant, décroît comme $(1-1/N)^n$.}\label{appF:fig:wf}
\end{figure}

\begin{exercise}[Exemple 4]\label{appF:exo4}
(a) Démontrer que $\lambda_2=1-1/N$ est valeur propre de $P$. On pourra appliquer $P$ au polynôme $i(N-i)$, après avoir
calculé $\mathbb E[X(N-X)]$ pour $X$ binomiale. (b) En déduire que $|P^n(i,N)-i/N|\le C(1-1/N)^n$. (c) Modifier le
script pour Moran au lieu de Wright--Fisher, et comparer les vitesses à $N$ égal.
\end{exercise}

% =====================================================================================================
\section{Exemple 5 -- Graphe isotherme contre étoile}\label{appF:ex5}

\paragraph{La question.} On place un mutant de valeur sélective $r$ sur un sommet d'un graphe, et la population évolue
par naissance--mort (le modèle de N4, chapitre~4). Le théorème isotherme \lean{ILENE} (\tierA) affirme ceci : si la
matrice de remplacement est doublement stochastique, la probabilité de fixation est celle de Moran,
$(1-r^{-1})/(1-r^{-n})$, \emph{quel que soit} le sommet de départ. Que se passe-t-il sur l'étoile, qui n'est pas
isotherme ?

\paragraph{Le code.} L'état est l'ensemble $S$ des mutants, codé par un masque de bits, ce qui donne $2^n$ états. Les
probabilités de fixation résolvent un système linéaire, que l'on résout en fractions par la méthode de Gauss :
\lstinputlisting[language=Python,style=appFcode,firstline=37,lastline=48,firstnumber=37]{code/examples/python/f05_graph_fixation.py}
\lstinputlisting[language=Python,style=appFcode,firstline=51,lastline=76,firstnumber=51]{code/examples/python/f05_graph_fixation.py}
\appFsortie{f05_graph_fixation.txt}

\paragraph{Lecture.} Sur le cycle à 4 sommets, qui est isotherme, la valeur est $8/15$ partout : c'est celle de Moran.
Sur l'étoile à 3 sommets, on retrouve les valeurs de N4, toutes trois \tierA{} :
\begin{itemize}
\item $5/12$ pour un mutant au centre ;
\item $2/3$ pour un mutant sur une feuille ;
\item $7/12$ pour un départ uniforme, contre $4/7$ pour Moran.
\end{itemize}
Les valeurs de l'étoile à 4 sommets, $392/1265$, $168/253$ et $728/1265$, ne sont que \tierX{} dans N4 : elles
sortent d'un script non vérifié, ici retrouvées par un second script. Les deux dernières lignes montrent pourquoi le
nombre de mutants n'est pas une chaîne de Markov sur le cycle. Avec deux mutants, la probabilité de monter d'un cran
vaut $1/3$ s'ils sont voisins et $2/3$ s'ils sont opposés. Mais le rapport $P^+/P^-$ vaut $r=2$ dans les deux cas :
c'est le biais par configuration que démontre N4.

\begin{figure}[ht]\centering
\includegraphics[width=\textwidth]{appF_graph_fixation.pdf}
\caption{Exemple~5, calcul exact. (a) Le cycle et l'étoile à 4 sommets. (b) Probabilité de fixation d'un mutant en
fonction de $r$. Le cycle suit Moran quel que soit le sommet. Sur l'étoile, un mutant au centre est désavantagé et un
mutant sur une feuille avantagé ; en moyenne, l'étoile amplifie la sélection pour $r>1$ dans cet échantillon de valeurs.}
\label{appF:fig:graph}
\end{figure}

\begin{exercise}[Exemple 5]\label{appF:exo5}
(a) Vérifier à la main que, sur l'étoile à 3 sommets avec $r=2$, l'état \texttt{001} (mutant au centre) passe à
\texttt{011} et à \texttt{101} avec la probabilité $1/4$ chacun, et à l'état vide avec la probabilité $1/2$. (b) Écrire
le graphe complet $K_4$ (avec $W(i,j)=1/3$ pour $i\ne j$) et prédire le résultat avant de lancer le programme. (c) Pour
l'étoile à $n$ sommets, conjecturer la limite de la probabilité moyenne quand $n\to\infty$, puis comparer à la littérature
(Lieberman, Hauert et Nowak, 2005). Une conjecture tirée de quelques valeurs exactes est de niveau~X.
\end{exercise}

% =====================================================================================================
\section{Exemple 6 -- Phonons et énergie de point zéro}\label{appF:ex6}

\paragraph{La question.} Au chapitre~3, la conjecture QU-1 demandait si le réseau triangulaire minimise aussi
l'énergie de point zéro harmonique $E_{zp}=\tfrac12\langle\omega_1+\omega_2\rangle_{\mathrm{ZB}}$. Pour la gaussienne
$e^{-\pi ar^2}$, le programme P2 a trouvé le contraire (\tierX). Sur $a\in(0{,}37668;1{,}58148)$, le réseau
triangulaire est un \emph{maximum} local de $E_{zp}$. Peut-on le retrouver avec trente lignes de numpy ?

\paragraph{Le code.} La matrice dynamique $D(\mathbf k)=\sum_{R\ne0}(1-\cos\mathbf k\cdot R)\,\mathrm{Hess}\,V(R)$ est
$2\times2$, donc ses valeurs propres ont une forme close. La moyenne sur la zone de Brillouin est prise sur une grille
$n\times n$ de la maille réciproque. La grille fractionnaire est la même pour toutes les formes $\tau$, ce qui fait
disparaître l'erreur de grille dans les différences :
\lstinputlisting[language=Python,style=appFcode,firstline=41,lastline=67,firstnumber=41]{code/examples/python/f06_phonons_zero_point.py}
\appFsortie{f06_phonons_zero_point.txt}

\paragraph{Lecture.} Les témoins positifs sont les valeurs de P2. Pour $a=1$, le script retrouve
$E_{zp}(\rho)=1{,}5033523522$ à $3\times10^{-8}$ près sur la grille $n=192$. Il retrouve aussi le contre-exemple de P2,
le losange $\tau=0{,}5+0{,}966025\,i$, plus bas de $1{,}3074\times10^{-2}$. La courbure $c(a)$ change de signe comme
annoncé :
\begin{itemize}
\item elle est négative pour $a=0{,}5$ et $a=1$, où $\rho$ est un maximum local ;
\item elle est positive pour $a=0{,}25$ et $a=2$, où $\rho$ est un minimum local.
\end{itemize}
Les quatre valeurs de P2 sont retrouvées au millième près, et la courbure ne bouge pas entre les grilles $48$ et $96$.
Tout cela reste \tierX : c'est un second calcul flottant, pas une preuve. Le certificat B2 du chapitre~3 traite le cas
$a=1$ avec une arithmétique de boules.

\begin{figure}[ht]\centering
\includegraphics[width=\textwidth]{appF_phonons.pdf}
\caption{Exemple~6 \tierX. (a) Les deux branches de phonons du réseau triangulaire pour $V=e^{-\pi r^2}$, le long de
$\Gamma$--M--K--$\Gamma$ ; elles s'annulent en $\Gamma$, comme pour tout cristal (modes acoustiques). (b) Variation relative de
$E_{zp}$ quand on déforme $\rho$ en $\rho+x$ : une parabole vers le bas pour $a=1$ (maximum local), vers le haut pour
$a=0{,}25$ et $a=2$.}\label{appF:fig:phonons}
\end{figure}

\begin{exercise}[Exemple 6]\label{appF:exo6}
(a) Démontrer que $D(\mathbf k)$ est symétrique et que $D(\mathbf 0)=0$. Pourquoi $\omega(\mathbf k)$ est-il linéaire en
$|\mathbf k|$ près de $\Gamma$ ? (b) Modifier \texttt{curvature} pour tracer $c(a)$ sur $a\in[0{,}3;1{,}7]$ et localiser
numériquement les deux changements de signe (P2 : $0{,}37668$ et $1{,}58148$). (c) Pourquoi l'énoncé \og le réseau
triangulaire est un maximum local de $E_{zp}$ pour $a=1$\fg{} ne peut-il pas devenir de niveau~A avec ce seul code ?
\end{exercise}

% =====================================================================================================
\section[Exemple 7 -- CVODE en Rust]{Exemple 7 -- CVODE en Rust : une EDO raide, puis l'équation de Kolmogorov}\label{appF:ex7}

\paragraph{La question.} rusty-SUNDIALS est une réécriture en Rust de SUNDIALS. Son solveur CVODE intègre les équations
différentielles raides par les formules BDF. Comment l'appeler depuis une crate indépendante ? Et pourquoi un schéma
explicite échoue-t-il là où BDF réussit ?

\paragraph{Une EDO raide.} Le problème de Prothero--Robinson, $y'=\lambda(y-\cos t)-\sin t$ avec $y(0)=2$ et
$\lambda=-10^4$, a pour solution exacte $y=\cos t+e^{\lambda t}$. Voici l'appel à CVODE (fonction de la bibliothèque de
la crate) :
\lstinputlisting[language=Rust,style=appFcode,firstline=36,lastline=68,firstnumber=36]{code/examples/rust/src/lib.rs}
\appFsortie{f07a_tb_stiff.txt}

\paragraph{Lecture.} CVODE suit la solution exacte à $6\times10^{-9}$ près pendant la phase transitoire, où $e^{\lambda t}$
s'éteint en $10^{-4}$, puis à $3\times10^{-11}$ près jusqu'à $t=10$. Il y parvient en 506 pas. Le schéma d'Euler explicite
n'est stable que si $|1+h\lambda|\le1$, c'est-à-dire $h\le2/|\lambda|=2\times10^{-4}$ :
\begin{itemize}
\item avec $h=10^{-4}$, il réussit, mais en $10\,000$ pas ;
\item avec $h=2{,}5\times10^{-4}$, l'erreur est multipliée par $|1+h\lambda|=1{,}5$ à chaque pas, atteint $4\times10^{17}$
en 100 pas et finit en \texttt{NaN}.
\end{itemize}
C'est le témoin négatif, et la définition même de la raideur.

\paragraph{L'équation de Kolmogorov de Moran.} La chaîne de Moran (chapitre~4) uniformisée en temps continu obéit à
$\dot\pi=\pi(P-I)$. La masse absorbée en $N$ tend vers la probabilité de fixation
$\varphi_i=(1-r^{-i})/(1-r^{-N})$. Le jacobien est la matrice constante $(P-I)^{\mathsf T}$. On double le temps de
sortie jusqu'à ce que la masse intérieure tombe sous $10^{-12}$ :
\lstinputlisting[language=Rust,style=appFcode,firstline=125,lastline=172,firstnumber=125]{code/examples/rust/src/lib.rs}
\appFsortie{f07b_tb_moran_forward.txt}

\paragraph{Lecture.} Pour $N=3$, $r=2$ et $i=1$, CVODE donne $4/7$ à $9{,}5\times10^{-14}$ près. Pour $N=10$, les 27
valeurs ($r=0{,}8$, $1$ et $1{,}5$, $i=1,\dots,9$) suivent la formule à $10^{-12}$ près. Les deux témoins négatifs sont
refusés :
\begin{itemize}
\item le noyau identité, qui n'absorbe jamais ;
\item la formule à exposant inversé, qui donne $1/7$.
\end{itemize}
La crate \texttt{qf-popgen} du chapitre~5 fait le même calcul avec un certificat d'absorption plus complet. Elle n'est
pas encore fusionnée dans \texttt{main}, c'est pourquoi cette crate réécrit l'équation directement.

\begin{figure}[ht]\centering
\includegraphics[width=\textwidth]{appF_rust_moran.pdf}
\caption{Exemple~7 \tierX, tracé depuis les fichiers CSV écrits par \texttt{tb\_moran\_forward}. (a) Probabilité de fixation pour
$N=10$ : CVODE (marqueurs) et formule exacte (traits). (b) $N=3$, $r=2$, $i=1$ : la masse en $N$ monte vers $4/7$ et la
masse intérieure s'éteint, aux temps de sortie $0{,}25\cdot2^k$.}\label{appF:fig:moran}
\end{figure}

\begin{exercise}[Exemple 7]\label{appF:exo7}
(a) Pour $y'=\lambda y$, montrer que le schéma d'Euler implicite (BDF d'ordre~1) donne $y_{n+1}=y_n/(1-h\lambda)$, stable
pour tout $h>0$ si $\lambda<0$. (b) Pourquoi la masse intérieure décroît-elle exponentiellement, et quel taux
attendre pour $N=3$, $r=2$ ? On calculera les valeurs propres de la restriction de $P-I$ aux états intérieurs.
(c) Ajouter un témoin : un noyau dont une ligne somme à $1{,}001$. Quelle quantité calculée par \texttt{absorb} le
détecte ?
\end{exercise}

% =====================================================================================================
\section{Exemple 8 -- Relaxer un petit cristal de Yukawa}\label{appF:ex8}

\paragraph{La question.} 64 particules interagissent par $e^{-2r}/r$ dans une boîte périodique. Si on les lâche d'une
position quelconque et qu'elles descendent le gradient d'énergie ($\dot x=F(x)$), retrouvent-elles le réseau
triangulaire ? Le théorème dit seulement que le réseau triangulaire est le \emph{minimum global} (\tierA,
conditionnel). Il ne dit pas que la descente l'atteint.

\paragraph{Le code.} La crate appelle l'API de \texttt{qf-yukawa-crystal} :
\begin{itemize}
\item \texttt{ExactBox} pour l'énergie, les forces et la hessienne sommées sur les images jusqu'à $R=12$ ;
\item \texttt{lattice\_energy\_exact} pour l'énergie exacte du réseau.
\end{itemize}
La relaxation elle-même est une intégration CVODE dont le jacobien est l'opposé de la hessienne (fonction
\texttt{relax\_traced} de la bibliothèque). Voici le programme principal :
\lstinputlisting[language=Rust,style=appFcode,firstline=60,lastline=105,firstnumber=60]{code/examples/rust/src/bin/tb_yukawa_relax.rs}
\appFsortie{f08_tb_yukawa_relax.txt}

\paragraph{Lecture.} Les trois départs se comportent comme prévu.
\begin{itemize}
\item \emph{Réseau triangulaire perturbé} (témoin positif) : il revient à l'énergie exacte du réseau, à $10^{-16}$ près.
Cette énergie, $0{,}393542848871227$ à $R=12$, diffère de $6\times10^{-11}$ de la valeur à $R=20$ enregistrée par
\texttt{exact\_refs.jsonl}, qui vaut $0{,}39354284893239$ ; c'est la queue omise.
\item \emph{Départ aléatoire} : il s'arrête $2{,}1\times10^{-3}$ \emph{au-dessus}, en un point stationnaire
(forces inférieures à $10^{-9}$) qui présente des défauts (figure~\ref{appF:fig:yukawa}b). Nous n'avons pas vérifié la
hessienne : rien ne dit ici qu'il s'agit d'un minimum local. Une descente de gradient ne trouve pas toujours le
minimum global.
\item \emph{Réseau carré dans une boîte carrée} (témoin négatif) : il descend sous l'énergie du carré, ce qui indique que
le réseau carré n'est pas un minimum, mais il reste au-dessus du triangulaire, que la boîte $8\times8$ ne peut pas
contenir.
\end{itemize}

\begin{figure}[ht]\centering
\includegraphics[width=\textwidth]{appF_rust_yukawa.pdf}
\caption{Exemple~8 \tierX. (a) Énergie par particule au-dessus du réseau triangulaire pendant la relaxation CVODE, aux
temps de sortie $0{,}01\cdot2^k$. (b) État final du départ aléatoire : des domaines triangulaires et des défauts. (c) État
final dans la boîte carrée.}\label{appF:fig:yukawa}
\end{figure}

\begin{exercise}[Exemple 8]\label{appF:exo8}
(a) Pourquoi la boîte $8a\times4\sqrt3\,a$ peut-elle contenir exactement le réseau triangulaire $8\times8$ de pas $a$,
alors que la boîte $8\times8$ ne le peut pas ? (b) Changer la graine du départ aléatoire (11) pour cinq autres valeurs.
Combien de départs atteignent le réseau triangulaire à $10^{-10}$ près ? (c) Ce résultat contredit-il le théorème ?
\end{exercise}

% =====================================================================================================
\section{Exemple 9 -- Les invariants du modèle en couches \texorpdfstring{$\sigma=3$}{sigma=3}}\label{appF:ex9}

\paragraph{La question.} Le modèle en couches de la section~\ref{ch5:U4} conserve deux quantités :
\begin{itemize}
\item $S=\sum|v_n|^2$ ;
\item $H_3$, dont la partie quadratique vaut $2D\Omega$, où $\Omega$ est l'enstrophie.
\end{itemize}
Lean en déduit la borne $2D\,\Omega(t)\le H_3(0)+S(0)^{3/2}$ le long des solutions (\tierAs, chapitre~5). Un
intégrateur numérique conserve-t-il ces invariants ? Et la borne tient-elle sur une trajectoire ?

\paragraph{Le code.} Six couches, $D=0{,}05$, intégration jusqu'à $T=2$ avec \texttt{integrate\_cvode} de
\texttt{qf-shell-cascade}, qui évalue les invariants à chaque pas interne de CVODE. Le témoin négatif multiplie un
couplage par $1+10^{-2}$ :
\lstinputlisting[language=Rust,style=appFcode,firstline=18,lastline=63,firstnumber=18]{code/examples/rust/src/bin/tb_shell_sigma3.rs}
\appFsortie{f09_tb_shell_sigma3.txt}

\paragraph{Lecture.} Pour le modèle exact, les invariants dérivent de $5\times10^{-10}$ pour $S$ et de $9\times10^{-8}$ pour
$H_3$, rapporté à l'échelle de ses termes. Ces dérives relèvent de la tolérance du solveur, pas du modèle. La borne de
Lean est loin d'être atteinte, avec un rapport maximal de $0{,}29$. Avec le couplage faussé, les dérives montent à
$3\times10^{-3}$ : la conservation tient à l'annulation exacte des termes non linéaires, que le témoin négatif brise.
Nous nous limitons à six couches : la rotation la plus rapide, $\omega_n=Dk_n^3$ avec $k_n=2^n$, est multipliée par~8 à
chaque couche ajoutée, et CVODE doit la résoudre pas à pas.

\begin{figure}[ht]\centering
\includegraphics[width=\textwidth]{appF_rust_shell.pdf}
\caption{Exemple~9 \tierX, tracé depuis le CSV écrit par \texttt{tb\_shell\_sigma3} (un pas interne sur vingt). (a) Dérive
relative des deux invariants. (b) Le rapport $2D\Omega(t)/(H_3(0)+S(0)^{3/2})$ reste sous~1, comme l'impose le théorème
de Lean.}\label{appF:fig:shell}
\end{figure}

\begin{exercise}[Exemple 9]\label{appF:exo9}
(a) Démontrer que $\frac{d}{dt}\sum|v_n|^2=0$ pour le modèle sans dissipation. On regroupera les termes de couplage
deux à deux. (b) Remplacer \texttt{sigma: 3.0} par \texttt{2.0}. Les invariants $S$ et $H_2$ sont-ils encore conservés ?
Pourquoi la borne d'enstrophie ne se déduit-elle plus de la même manière ?
\end{exercise}

% =====================================================================================================
\section[Exemple 10 -- LeanFlow]{Exemple 10 -- LeanFlow : certifier des fractions, rejeter un mensonge}\label{appF:ex10}

\paragraph{La question.} L'exemple~\ref{appF:ex5} a calculé en Python la table exacte des probabilités de fixation de
l'étoile à 3 sommets. Python peut se tromper, ou être mal programmé. Comment faire vérifier ces fractions par le noyau
de Lean, sans écrire une ligne de Lean ?

\paragraph{Le code.} \texttt{lean.certify} de LeanFlow reçoit des affirmations sur des rationnels. Il les traduit en
énoncés sur les entiers, les fait décider par \lean{decide +kernel} (tactique \texttt{decide\_kernel}), puis audite
\texttt{\#print axioms}. Un oracle exact en Python doit rendre le même verdict que le noyau, sinon le résultat est une
erreur, jamais un succès. Voici la construction de l'affirmation « le tableau $x_S$ est harmonique » :
\lstinputlisting[language=Python,style=appFcode,firstline=34,lastline=49,firstnumber=34]{code/examples/python/f10_leanflow_certify.py}
\appFsortie{f10_leanflow_certify.txt}
Voici le plus court des deux certificats engendrés, tel que le noyau l'a décidé :
\lstinputlisting[language=Lean,style=appFcode]{code/examples/lean/f10_generated_PosMeanAboveMoran.lean}

\paragraph{Lecture.} Les deux affirmations vraies sont acceptées :
\begin{itemize}
\item les six équations du premier pas, satisfaites exactement par le tableau $(0,\tfrac5{12},\tfrac23,\tfrac56,
\tfrac23,\tfrac56,\tfrac{14}{15},1)$ ;
\item l'inégalité $7/12>4/7$.
\end{itemize}
Les deux affirmations plantées sont rejetées, avec \lean{sorryAx} dans la liste des axiomes :
\begin{itemize}
\item le vecteur de Moran $(0,\tfrac47,\tfrac67,1)$ selon $|S|$ viole les six équations intérieures, mais il satisfait
les deux conditions au bord (\texttt{cert\_0} et \texttt{cert\_1} ne sont pas refusés) ;
\item l'inégalité renversée $7/12<4/7$ est refusée.
\end{itemize}
LeanFlow ne contrôle que \texttt{\#print axioms}, qui peut être trompé (chapitre~2). Le script complète donc la porte du
niveau~A. Chaque source acceptée est recompilée en \texttt{.olean} avec Lean~4.34.1 puis rejouée par
\texttt{leanchecker}, avec le code de retour~0 pour les deux. Aucun des mots \texttt{skipKernelTC}, \texttt{addDecl},
\texttt{native\_decide} ou \texttt{implemented\_by} n'y figure. Le témoin du rejeu, un module inexistant, sort avec le
code~1. Les deux certificats sont donc \tierA, mais pour l'\emph{énoncé arithmétique seulement}. Que ce tableau soit la
probabilité de fixation de l'étoile demande en plus l'unicité de la solution harmonique (chapitre~4) et l'identification
des transitions de Python à celles de \lean{bdK}. Ce raccord est un argument humain \tierC.

\begin{exercise}[Exemple 10, LeanFlow]\label{appF:exo10}
(a) Certifier de la même manière le tableau du cycle à 4 sommets ($2^4$ états, valeur $8/15$ pour tout mutant
seul). (b) Planter une erreur d'un seul ulp : remplacer $5/12$ par la fraction du flottant \texttt{float(5/12)}.
LeanFlow l'accepte-t-il ? Pourquoi un certificat sur des rationnels est-il plus exigeant qu'une comparaison
flottante à $10^{-15}$ près ?
\end{exercise}

% =====================================================================================================
\section{Exemple 11 -- Lire et vérifier le théorème ILENE}\label{appF:ex11}

\paragraph{La question.} Le chapitre~5 cite \lean{ILENE}, la version à un mutant du théorème isotherme. Comment le
lecteur peut-il voir, sur sa propre machine, l'énoncé exact que Lean a vérifié et les axiomes dont dépend sa preuve ?
Et comment distinguer une vraie preuve d'une fausse ?

\paragraph{Le code.} Un fichier de 21 lignes importe le module compilé et pose quatre questions à Lean :
\lstinputlisting[language=Lean,style=appFcode]{code/examples/lean/f11_ilene.lean}
On le compile avec Lean~4.34.1 et Mathlib \texttt{d13f23b7}, par le script \texttt{lean\_pinned.py} du
dossier \texttt{scripts/certified\_numerics/}. Il donne :
\appFsortie{f11_ilene.txt}

\paragraph{Lecture.} La sortie se lit en quatre temps.
\begin{itemize}
\item \texttt{\#check} affiche l'énoncé complet. Les hypothèses sont : $0<n$, $0<r$, $r\ne1$, $W$ isotherme et fortement
connexe. La conclusion est une limite, celle de $P^t(\{v\},V)$ vers $(1-r^{-1})/(1-r^{-n})$, pour tout sommet $v$.
\item Les axiomes sont exactement ceux de la liste blanche. Le module des acronymes et ses sources ont été rejoués
par \texttt{leanchecker}, avec le code de retour~0 (enregistrement \texttt{acronymes\_leanchecker.json}) :
\lean{ILENE} est donc \tierA.
\item L'\texttt{example} vérifie, par \lean{norm_num}, que la valeur limite vaut $8/15$ pour $n=4$ et $r=2$. C'est
celle du cycle de l'exemple~\ref{appF:ex5}.
\item Le témoin planté affirme une valeur fausse, $4/7$, « prouvée » par \lean{sorry}. Le fichier compile quand même,
avec le code de retour~0 et un simple avertissement. Seule la ligne \texttt{\#print axioms} trahit la fraude : elle
affiche \lean{sorryAx}.
\end{itemize}
C'est la leçon du chapitre~2 : un code de retour nul ne prouve rien, il faut lire les axiomes et rejouer le noyau.

\begin{exercise}[Exemple 11, Lean]\label{appF:exo11}
(a) Ajouter \texttt{\#check @KAL} et \texttt{\#print axioms AURELIA}, et relire les énoncés à l'aide du tableau des
acronymes (section~\ref{ch5:acronymes}). (b) Remplacer \lean{by sorry} par \lean{by norm_num} dans le témoin
planté. Que répond Lean, et avec quel code de retour ? (c) Pourquoi l'égalité $1/2=2^{-1}$, utilisée implicitement en
comparant l'\texttt{example} à l'énoncé de \lean{ILENE}, est-elle un raccord à vérifier, et comment le faire en Lean ?
\end{exercise}

% =====================================================================================================
\section{Exemple 12 -- Contrôle croisé : CVODE contre SciPy}\label{appF:ex12}

\paragraph{La question.} Deux codes écrits par la même équipe peuvent partager la même erreur. Les nombres de
l'exemple~\ref{appF:ex7} sont-ils confirmés par des méthodes \emph{indépendantes} ?

\paragraph{Le code.} Le script relit les sorties du binaire Rust et recalcule les mêmes quantités de deux façons :
\begin{itemize}
\item pour le problème raide, avec la méthode Radau de SciPy, une méthode de Runge--Kutta implicite et non BDF ;
\item pour Moran, avec l'exponentielle de matrice $\pi(t)=e_i\,e^{t(P-I)}$, sans aucun pas de temps.
\end{itemize}
\lstinputlisting[language=Python,style=appFcode,firstline=56,lastline=78,firstnumber=56]{code/examples/python/f12_scipy_crosscheck.py}
\appFsortie{f12_scipy_crosscheck.txt}

\paragraph{Lecture.} Après la phase transitoire, CVODE et Radau s'accordent à $3\times10^{-11}$ près. Aux deux premiers
instants, l'écart de $5\times10^{-9}$ est celui de CVODE à la solution exacte : la valeur de Radau coïncide avec
$\cos t+e^{\lambda t}$, et la tolérance de CVODE était $10^{-8}$. Pour Moran, CVODE et l'exponentielle de matrice
s'accordent à $1{,}3\times10^{-11}$ près aux dix temps de sortie. Le témoin négatif décale un temps de $10^{-3}$, et
l'écart, de $8\times10^{-4}$, devient immédiatement visible : la comparaison a donc le pouvoir de détecter une erreur.

\begin{exercise}[Exemple 12]\label{appF:exo12}
(a) Remplacer \texttt{method="Radau"} par \texttt{"RK45"}, une méthode explicite. Combien d'évaluations faut-il, et
pourquoi ? (b) Pour Moran, l'exponentielle de matrice est exacte à l'arrondi près. Pourquoi ne remplace-t-elle pas
CVODE pour $N=1000$ ?
\end{exercise}

% =====================================================================================================
\section{Bilan}\label{appF:bilan}

Les douze exemples ont été exécutés le 2026-10-26, et tous leurs témoins se sont comportés comme prévu. Ils illustrent
les trois manières de traiter un nombre dans ce livre.
\begin{itemize}
\item \emph{Le calcul flottant} (exemples~1, 2, 3, 6 à 9 et 12) explore, contrôle et réfute. Il retrouve les valeurs
publiées à la précision attendue, mais il reste de niveau~X.
\item \emph{Le calcul exact en rationnels} (exemples~4 et~5) supprime l'arrondi, mais pas le risque d'une erreur de
programmation.
\item \emph{Le noyau de Lean} (exemples~10 et~11) supprime ce risque pour l'énoncé vérifié, et pour lui seul. Le lien
entre cet énoncé et la question posée reste un raccord, à vérifier à part.
\end{itemize}
